%==============================================================================
%  Referencias (formato APA, 7.a edición)
%  Cada \bibitem lleva entre corchetes la etiqueta que aparece en el texto:
%    \bibitem[Autor(Año)]{clave}           -> (Autor, Año)
%    \bibitem[{Autor1 \& Autor2}(Año)]{..}  -> (Autor1 & Autor2, Año)
%    \bibitem[{Autor1 et~al.}(Año)]{..}     -> (Autor1 et al., Año)
%  Mantener la lista en orden alfabético.
%==============================================================================
\cleardoublepage
\phantomsection
\addcontentsline{toc}{chapter}{\bibname}

\begin{thebibliography}{99}

\bibitem[{Adams et~al.}(1988)]{adams1988}
Adams, J., Balas, E., \& Zawack, D. (1988). The shifting bottleneck procedure
for job shop scheduling. \textit{Management Science}, \textit{34}(3), 391--401.

\bibitem[{Amadini et~al.}(2014a)]{amadini2014mzn2feat}
Amadini, R., Gabbrielli, M., \& Mauro, J. (2014a). An enhanced features
extractor for a portfolio of constraint solvers. En \textit{Proceedings of the
29th Annual ACM Symposium on Applied Computing (SAC 2014)} (pp. 1357--1359).
ACM.

\bibitem[{Amadini et~al.}(2014b)]{amadini2014sunny}
Amadini, R., Gabbrielli, M., \& Mauro, J. (2014b). SUNNY: A lazy portfolio
approach for constraint solving. \textit{Theory and Practice of Logic
Programming}, \textit{14}(4--5), 509--524.
\url{https://doi.org/10.1017/S1471068414000179}

\bibitem[{Amadini et~al.}(2015)]{amadini2015sunnycp}
Amadini, R., Gabbrielli, M., \& Mauro, J. (2015). SUNNY-CP: A sequential CP
portfolio solver. En \textit{Proceedings of the 30th Annual ACM Symposium on
Applied Computing (SAC 2015)}. ACM. \url{https://doi.org/10.1145/2695664.2695741}

\bibitem[{Applegate \& Cook}(1991)]{applegate1991}
Applegate, D., \& Cook, W. (1991). A computational study of the job-shop
scheduling problem. \textit{ORSA Journal on Computing}, \textit{3}(2),
149--156.

\bibitem[{Beasley}(s.~f.)]{orlibjobshop}
Beasley, J. E. (s. f.). \textit{OR-Library: Job shop scheduling} [Conjunto de
datos, archivo jobshop1.txt]. Recuperado el 10 de noviembre de 2026, de
\url{https://people.brunel.ac.uk/~mastjjb/jeb/orlib/jobshopinfo.html}

\bibitem[Beasley(1990)]{beasley1990}
Beasley, J. E. (1990). OR-Library: Distributing test problems by electronic
mail. \textit{Journal of the Operational Research Society}, \textit{41}(11),
1069--1072.

\bibitem[{Bengio et~al.}(2021)]{bengio2021}
Bengio, Y., Lodi, A., \& Prouvost, A. (2021). Machine learning for
combinatorial optimization: A methodological tour d'horizon. \textit{European
Journal of Operational Research}, \textit{290}(2), 405--421.

\bibitem[{Bischl et~al.}(2016)]{bischl2016aslib}
Bischl, B., Kerschke, P., Kotthoff, L., Lindauer, M., Malitsky, Y., Fréchette,
A., Hoos, H., Hutter, F., Leyton-Brown, K., Tierney, K., \& Vanschoren, J.
(2016). ASlib: A benchmark library for algorithm selection.
\textit{Artificial Intelligence}, \textit{237}, 41--58.

\bibitem[Breiman(2001)]{breiman2001}
Breiman, L. (2001). Random forests. \textit{Machine Learning}, \textit{45}(1),
5--32.

\bibitem[{Brucker et~al.}(1994)]{brucker1994}
Brucker, P., Jurisch, B., \& Sievers, B. (1994). A branch and bound algorithm
for the job-shop scheduling problem. \textit{Discrete Applied Mathematics},
\textit{49}(1--3), 107--127.

\bibitem[{Carlier \& Pinson}(1989)]{carlier1989}
Carlier, J., \& Pinson, E. (1989). An algorithm for solving the job-shop
problem. \textit{Management Science}, \textit{35}(2), 164--176.

\bibitem[{Carlier \& Pinson}(1994)]{carlier1994}
Carlier, J., \& Pinson, E. (1994). Adjustment of heads and tails for the
job-shop problem. \textit{European Journal of Operational Research},
\textit{78}(2), 146--161.

\bibitem[{Carrera-Rivera et~al.}(2022)]{carrera2022}
Carrera-Rivera, A., Ochoa, W., Larrinaga, F., \& Lasa, G. (2022). How-to
conduct a systematic literature review: A quick guide for computer science
research. \textit{MethodsX}, \textit{9}, 101895.
\url{https://doi.org/10.1016/j.mex.2022.101895}

\bibitem[Chu(2011)]{chu2011}
Chu, G. (2011). \textit{Improving combinatorial optimization} [Tesis
doctoral]. University of Melbourne.

\bibitem[COSEAL(2024)]{coseal2024}
COSEAL. (2024). \textit{COSEAL: Configuration and Selection of Algorithms}.
\url{https://www.coseal.net/}

\bibitem[{Cover \& Hart}(1967)]{cover1967}
Cover, T., \& Hart, P. (1967). Nearest neighbor pattern classification.
\textit{IEEE Transactions on Information Theory}, \textit{13}(1), 21--27.

\bibitem[{Demirkol et~al.}(1998)]{demirkol1998}
Demirkol, E., Mehta, S., \& Uzsoy, R. (1998). Benchmarks for shop scheduling
problems. \textit{European Journal of Operational Research}, \textit{109}(1),
137--141.

\bibitem[{Fisher \& Thompson}(1963)]{fisher1963}
Fisher, H., \& Thompson, G. L. (1963). Probabilistic learning combinations of
local job-shop scheduling rules. En J. F. Muth \& G. L. Thompson (Eds.),
\textit{Industrial scheduling} (pp. 225--251). Prentice-Hall.

\bibitem[Friedman(2001)]{friedman2001}
Friedman, J. H. (2001). Greedy function approximation: A gradient boosting
machine. \textit{The Annals of Statistics}, \textit{29}(5), 1189--1232.

\bibitem[{Garey et~al.}(1976)]{garey1976}
Garey, M. R., Johnson, D. S., \& Sethi, R. (1976). The complexity of flowshop
and jobshop scheduling. \textit{Mathematics of Operations Research},
\textit{1}(2), 117--129.

\bibitem[{Gecode Team}(2019)]{gecode2019}
Gecode Team. (2019). \textit{Gecode: Generic constraint development
environment} (versión 6) [Software]. \url{https://www.gecode.org}

\bibitem[{Giffler \& Thompson}(1960)]{giffler1960}
Giffler, B., \& Thompson, G. L. (1960). Algorithms for solving
production-scheduling problems. \textit{Operations Research}, \textit{8}(4),
487--503.

\bibitem[{Gomes \& Selman}(2001)]{gomes2001}
Gomes, C. P., \& Selman, B. (2001). Algorithm portfolios. \textit{Artificial
Intelligence}, \textit{126}(1--2), 43--62.

\bibitem[Google(2025)]{perron2025}
Google. (2025). \textit{OR-Tools: CP-SAT solver} (versión 9.12) [Software].
\url{https://developers.google.com/optimization/cp/}

\bibitem[{Hoos et~al.}(2015)]{hoos2015aspeed}
Hoos, H., Kaminski, R., Lindauer, M., \& Schaub, T. (2015). aspeed: Solver
scheduling via answer set programming. \textit{Theory and Practice of Logic
Programming}, \textit{15}(1), 117--142.
\url{https://doi.org/10.1017/S1471068414000015}

\bibitem[{Hutter et~al.}(2014)]{hutter2014}
Hutter, F., Xu, L., Hoos, H. H., \& Leyton-Brown, K. (2014). Algorithm runtime
prediction: Methods \& evaluation. \textit{Artificial Intelligence},
\textit{206}, 79--111.

\bibitem[{Kerschke et~al.}(2019)]{kerschke2019}
Kerschke, P., Hoos, H. H., Neumann, F., \& Trautmann, H. (2019). Automated
algorithm selection: Survey and perspectives. \textit{Evolutionary
Computation}, \textit{27}(1), 3--45.

\bibitem[Kotthoff(2014)]{kotthoff2014}
Kotthoff, L. (2014). Algorithm selection for combinatorial search problems: A
survey. \textit{AI Magazine}, \textit{35}(3), 48--60.

\bibitem[{Ku \& Beck}(2016)]{ku2016}
Ku, W.-Y., \& Beck, J. C. (2016). Mixed integer programming models for job
shop scheduling: A computational analysis. \textit{Computers \& Operations
Research}, \textit{73}, 165--173.

\bibitem[Lawrence(1984)]{lawrence1984}
Lawrence, S. (1984). \textit{Resource constrained project scheduling: An
experimental investigation of heuristic scheduling techniques (Supplement)}.
Graduate School of Industrial Administration, Carnegie Mellon University.

\bibitem[{Lindauer et~al.}(2019)]{lindauer2019}
Lindauer, M., van Rijn, J. N., \& Kotthoff, L. (2019). The algorithm
selection competitions 2015 and 2017. \textit{Artificial Intelligence},
\textit{272}, 86--100. \url{https://doi.org/10.1016/j.artint.2018.10.004}

\bibitem[{Mirshekarian \& Šormaz}(2016)]{mirshekarian2016}
Mirshekarian, S., \& Šormaz, D. N. (2016). Correlation of job-shop scheduling
problem features with scheduling efficiency. \textit{Expert Systems with
Applications}, \textit{62}, 131--147.

\bibitem[{Nethercote et~al.}(2007)]{nethercote2007}
Nethercote, N., Stuckey, P. J., Becket, R., Brand, S., Duck, G. J., \& Tack,
G. (2007). MiniZinc: Towards a standard CP modelling language. En C. Bessière
(Ed.), \textit{Principles and Practice of Constraint Programming -- CP 2007}
(pp. 529--543). Springer.

\bibitem[{Nowicki \& Smutnicki}(1996)]{nowicki1996}
Nowicki, E., \& Smutnicki, C. (1996). A fast taboo search algorithm for the
job shop problem. \textit{Management Science}, \textit{42}(6), 797--813.

\bibitem[{Ohrimenko et~al.}(2009)]{ohrimenko2009}
Ohrimenko, O., Stuckey, P. J., \& Codish, M. (2009). Propagation via lazy
clause generation. \textit{Constraints}, \textit{14}(3), 357--391.

\bibitem[{O'Mahony et~al.}(2008)]{omahony2008cphydra}
O'Mahony, E., Hebrard, E., Holland, A., Nugent, C., \& O'Sullivan, B. (2008).
Using case-based reasoning in an algorithm portfolio for constraint solving.
En \textit{Proceedings of the 19th Irish Conference on Artificial Intelligence
and Cognitive Science (AICS 2008)}.

\bibitem[{Panwalkar \& Iskander}(1977)]{panwalkar1977}
Panwalkar, S. S., \& Iskander, W. (1977). A survey of scheduling rules.
\textit{Operations Research}, \textit{25}(1), 45--61.

\bibitem[{Pedregosa et~al.}(2011)]{pedregosa2011}
Pedregosa, F., Varoquaux, G., Gramfort, A., Michel, V., Thirion, B., Grisel,
O., Blondel, M., Prettenhofer, P., Weiss, R., Dubourg, V., Vanderplas, J.,
Passos, A., Cournapeau, D., Brucher, M., Perrot, M., \& Duchesnay, É. (2011).
Scikit-learn: Machine learning in Python. \textit{Journal of Machine Learning
Research}, \textit{12}, 2825--2830.

\bibitem[{Perron et~al.}(2023)]{perron2023}
Perron, L., Didier, F., \& Gay, S. (2023). The CP-SAT-LP solver (invited
talk). En R. H. C. Yap (Ed.), \textit{29th International Conference on
Principles and Practice of Constraint Programming (CP 2023)} (LIPIcs, Vol.
280, pp. 3:1--3:2). Schloss Dagstuhl -- Leibniz-Zentrum für Informatik.
\url{https://doi.org/10.4230/LIPIcs.CP.2023.3}

\bibitem[Pinedo(2022)]{pinedo2022}
Pinedo, M. L. (2022). \textit{Scheduling: Theory, algorithms, and systems}
(6.a ed.). Springer.

\bibitem[{Reijnen et~al.}(2023)]{reijnen2023}
Reijnen, R., Smit, I. G., Zhang, H., Wu, Y., Bukhsh, Z., \& Zhang, Y. (2023).
\textit{Job shop scheduling benchmark: Environments and instances for learning
and non-learning methods} (arXiv:2308.12794). arXiv.
\url{https://arxiv.org/abs/2308.12794}

\bibitem[Rice(1976)]{rice1976}
Rice, J. R. (1976). The algorithm selection problem. \textit{Advances in
Computers}, \textit{15}, 65--118.

\bibitem[{Rossi et~al.}(2006)]{rossi2006}
Rossi, F., van Beek, P., \& Walsh, T. (Eds.). (2006). \textit{Handbook of
constraint programming}. Elsevier.

\bibitem[{Schutt et~al.}(2011)]{schutt2011}
Schutt, A., Feydy, T., Stuckey, P. J., \& Wallace, M. G. (2011). Explaining
the cumulative propagator. \textit{Constraints}, \textit{16}(3), 250--282.

\bibitem[{Smit et~al.}(2025)]{smit2024gnn}
Smit, I. G., Zhou, J., Reijnen, R., Wu, Y., Chen, J., Zhang, C., Bukhsh, Z.,
Zhang, Y., \& Nuijten, W. (2025). Graph neural networks for job shop
scheduling problems: A survey. \textit{Computers \& Operations Research},
\textit{176}, 106914. \url{https://doi.org/10.1016/j.cor.2024.106914}

\bibitem[Smith-Miles(2008)]{smithmiles2009}
Smith-Miles, K. (2008). Cross-disciplinary perspectives on meta-learning for
algorithm selection. \textit{ACM Computing Surveys}, \textit{41}(1), Artículo 6.
\url{https://doi.org/10.1145/1456650.1456656}

\bibitem[{Storer et~al.}(1992)]{storer1992}
Storer, R. H., Wu, S. D., \& Vaccari, R. (1992). New search spaces for
sequencing problems with application to job shop scheduling.
\textit{Management Science}, \textit{38}(10), 1495--1509.

\bibitem[{Strassl \& Musliu}(2022)]{strassl2022}
Strassl, S., \& Musliu, N. (2022). Instance space analysis and algorithm
selection for the job shop scheduling problem. \textit{Computers \&
Operations Research}, \textit{141}, 105661.

\bibitem[{Streeter \& Smith}(2008)]{streeter2008}
Streeter, M., \& Smith, S. F. (2008). New techniques for algorithm portfolio
design. En \textit{Proceedings of the 24th Conference on Uncertainty in
Artificial Intelligence (UAI 2008)} (pp. 519--527). AUAI Press.

\bibitem[{Stuckey et~al.}(2014)]{stuckey2014}
Stuckey, P. J., Feydy, T., Schutt, A., Tack, G., \& Fischer, J. (2014). The
MiniZinc Challenge 2008--2013. \textit{AI Magazine}, \textit{35}(2), 55--60.

\bibitem[Taillard(1993)]{taillard1993}
Taillard, E. (1993). Benchmarks for basic scheduling problems.
\textit{European Journal of Operational Research}, \textit{64}(2), 278--285.

\bibitem[Weise(2019)]{weise2019jssp}
Weise, T. (2019). \textit{jsspInstancesAndResults: Results, data, and
instances of the job shop scheduling problem} [Repositorio de software].
GitHub. \url{https://github.com/thomasWeise/jsspInstancesAndResults}

\bibitem[{Wolpert \& Macready}(1997)]{wolpert1997}
Wolpert, D. H., \& Macready, W. G. (1997). No free lunch theorems for
optimization. \textit{IEEE Transactions on Evolutionary Computation},
\textit{1}(1), 67--82.

\bibitem[{Xu et~al.}(2010)]{xu2010hydra}
Xu, L., Hoos, H. H., \& Leyton-Brown, K. (2010). Hydra: Automatically
configuring algorithms for portfolio-based selection. En \textit{Proceedings of
the AAAI Conference on Artificial Intelligence}, \textit{24}(1), 210--216.

\bibitem[{Xu et~al.}(2008)]{xu2008satzilla}
Xu, L., Hutter, F., Hoos, H. H., \& Leyton-Brown, K. (2008). SATzilla:
Portfolio-based algorithm selection for SAT. \textit{Journal of Artificial
Intelligence Research}, \textit{32}, 565--606.

\bibitem[{Xu et~al.}(2025)]{xu2025learn}
Xu, M., Mei, Y., Zhang, F., \& Zhang, M. (2025). Learn to optimise for job
shop scheduling: A survey with comparison between genetic programming and
reinforcement learning. \textit{Artificial Intelligence Review}, \textit{58},
160. \url{https://doi.org/10.1007/s10462-024-11059-9}

\bibitem[{Yamada \& Nakano}(1992)]{yamada1992}
Yamada, T., \& Nakano, R. (1992). A genetic algorithm applicable to
large-scale job-shop problems. En R. Männer \& B. Manderick (Eds.),
\textit{Parallel Problem Solving from Nature 2} (pp. 281--290). Elsevier.

\end{thebibliography}
